% GAL1897_W08 — English translation of physical PDF pages 054–061.
% Primary authority: the frozen 1897 French diplomatic transcription from
% JP2 leaves L0053–L0060, navigated against the canonical 96-page 1897 PDF.
% The logical printed folio 25 is suppressed on the opening page; visible folios
% 26–32 remain recorded at their exact physical-page boundaries.
% Definite printed errors and unresolved witness variants are translated without
% silent repair and identified in adjacent notes from the frozen critical baseline.

\providecommand{\GalSourcePageStart}[4]{}%
\providecommand{\GalSourcePageBoundary}[4]{\space}%
\providecommand{\GalPrintedError}[2]{#1}%
\providecommand{\GalSourceError}[2]{#1}%
\providecommand{\GalWitnessVariant}[2]{#1}%
\providecommand{\GalHistoricalFootnoteText}[2]{\footnotetext{#2}}%
\providecommand{\GalHistoricalFootnote}[2]{\footnote{#2}}%
\providecommand{\GalArticleEndRule}{\par\smallskip\begin{center}\rule{0.18\linewidth}{0.4pt}\end{center}}%

\GalSourcePageStart{054}{L0053}{25}{suppressed}

% ENSEG:W08:PDF054:0001
\begin{center}
\rule{0.72\linewidth}{0.4pt}\par
\vspace{2.2em}
{\Large\bfseries II. --- POSTHUMOUS WORKS.\par}
\vspace{1.1em}
\rule{0.12\linewidth}{0.4pt}\par
\vspace{1.9em}
{\large\bfseries LETTER TO AUGUSTE CHEVALIER\textsuperscript{(*)}.\par}
\vspace{1.0em}
\rule{0.12\linewidth}{0.4pt}
\end{center}

% ENSEG:W08:PDF054:0002
\GalHistoricalFootnoteText{*}{Written on the eve of the author's death. (Published in 1832 in the
\emph{Revue encyclopédique}, September issue, page~568.) \textsc{(J. Liouville.)}}

% ENSEG:W08:PDF054:0003
\begin{flushleft}
\hspace{3.1em}My dear friend,
\end{flushleft}

% ENSEG:W08:PDF054:0004
I have done several new things in analysis.

% ENSEG:W08:PDF054:0005
Some concern the theory of equations; others, integral functions.

% ENSEG:W08:PDF054:0006
In the theory of equations, I investigated the cases in which equations
are solvable by radicals. This gave me occasion to study the theory more
deeply and to describe all possible transformations of an equation,
even when it is not solvable by radicals.

% ENSEG:W08:PDF054:0007
Three Memoirs can be made from all this.

% ENSEG:W08:PDF054:0008
The first is written, and, despite what Poisson said of it, I stand by
it, with the corrections I have made.

\medskip

% ENSEG:W08:PDF054:0009
The second contains some rather curious applications of the theory of
equations. Here is a summary of the most important things:

% ENSEG:W08:PDF054:0010
1\textsuperscript{o} From Propositions II and III of the first Memoir,
one sees a great difference between adjoining to an equation one root
of an auxiliary equation and adjoining all its roots.

% ENSEG:W08:PDF054:0011
In both cases, upon adjunction the group of the equation splits into
groups such that one passes from one to another by the same
substitution; but the condition that these groups have the same
substitutions certainly holds only in the second case. This is called
\emph{proper decomposition}.

% ENSEG:W08:PDF054:0012
In other words, when a group $G$ contains another group $H$,

\GalSourcePageBoundary{055}{L0054}{26}{26}

% ENSEG:W08:PDF055:0013
the group $G$ can be split into groups, each obtained by applying the
same substitution to the permutations of $H$, so that
% ENSEG:W08:PDF055:0014
\[
G=H+HS+HS'+\ldots.
\]
% ENSEG:W08:PDF055:0015
And it can also be divided into groups that all have the same
substitutions, so that
% ENSEG:W08:PDF055:0016
\[
G=H+TH+T'H+\ldots.
\]

% ENSEG:W08:PDF055:0017
These two kinds of decomposition do not ordinarily coincide. When they
do coincide, the decomposition is called \emph{proper}.

% ENSEG:W08:PDF055:0018
It is easy to see that, when the group of an equation admits no proper
decomposition, no matter how the equation is transformed, the groups
of the transformed equations will always have the same number of
permutations.

% ENSEG:W08:PDF055:0019
By contrast, when the group of an equation admits a proper
decomposition, so that it splits into $M$ groups of $N$ permutations,
one can solve the given equation by means of two equations: one will
have a group of $M$ permutations, the other one of $N$ permutations.

% ENSEG:W08:PDF055:0020
Thus, when one has exhausted on the group of an equation all possible
proper decompositions of that group, one arrives at groups that can be
transformed but whose permutations will always remain equal in number.

% ENSEG:W08:PDF055:0021
If these groups each have a prime number of permutations, the equation
will be solvable by radicals; otherwise, it will not.

% ENSEG:W08:PDF055:0022
The smallest number of permutations that an indecomposable group can
have, when this number is not prime, is $5.4.3$.

% ENSEG:W08:PDF055:0023
2\textsuperscript{o} The simplest decompositions are those obtained by
M.~Gauss's method.

% ENSEG:W08:PDF055:0024
Since these decompositions are evident even in the present form of the
group of the equation, it is unnecessary to dwell long on this subject.

% ENSEG:W08:PDF055:0025
What decompositions can be carried out on an equation that is not
simplified by M.~Gauss's method?

% ENSEG:W08:PDF055:0026
I have called \emph{primitive} the equations that cannot be simplified
by M.~Gauss's method; not that these equations are

\GalSourcePageBoundary{056}{L0055}{27}{27}

% ENSEG:W08:PDF056:0027
actually indecomposable, since they may even be solvable by radicals.

% ENSEG:W08:PDF056:0028
As a lemma for the theory of primitive equations solvable by radicals,
I published in June 1830, in the
\GalWitnessVariant{\emph{Bulletin de Férussac}}{W08-WV003}, an analysis
of the imaginaries in the theory of numbers.

% ENSEG:W08:PDF056:0029
In the attached material
% ENSEG:W08:PDF056:0125
\GalHistoricalFootnote{*}{Galois is speaking of the hitherto unpublished manuscripts that we are
publishing. \hfill \textsc{(J. Liouville.)}}
one will find the proof of the following theorems:

% ENSEG:W08:PDF056:0030
1\textsuperscript{o} For a primitive equation to be solvable by
radicals, it must have degree $\GalWitnessVariant{p^{\nu}}{W08-WV002}$,
where $p$ is prime.

\GalCriticalNote{W08-WV002}{Liouville's change of indices}{The 1897 edition
uses the indices $p,\nu$. Tannery's 1908 collation reports $p,n$ in the
manuscript and attributes the change to Liouville. The translation retains the
1897 notation; the manuscript form is recorded here only as witness history.}

\GalCriticalNote{W08-WV003}{Spelling of \emph{Férussac}}{The 1897 edition
prints \emph{Bulletin de Férussac}; Tannery's 1908 collation reports a
lowercase, unaccented form in the manuscript. The translation preserves the
bibliographic form printed in 1897 rather than importing the manuscript
spelling.}

% ENSEG:W08:PDF056:0031
2\textsuperscript{o} All the permutations of such an equation are of
the form
% ENSEG:W08:PDF056:0032
\[
x_{k,l,m,\ldots}\ \big|\
x_{ak+bl+cm+\ldots+h,\;a'k+b'l+c'm+\ldots+h',\;a''k+\ldots,\ldots},
\]
% ENSEG:W08:PDF056:0033
where $k$, $l$, $m$, \ldots\ are $\nu$ indices which, each taking $p$
values, designate all the roots. The indices are taken modulo $p$; that
is, the root will be the same when a multiple of $p$ is added to one of
the indices.

% ENSEG:W08:PDF056:0034
The group obtained by performing all substitutions of this linear form
contains in all
% ENSEG:W08:PDF056:0035
\[
p^{\nu}(p^{\nu}-1)(p^{\nu}-p)\ldots(p^{\nu}-p^{\nu-1})\text{ permutations.}
\]

% ENSEG:W08:PDF056:0036
It is far from being the case that, in this generality, the corresponding
equations are solvable by radicals.

% ENSEG:W08:PDF056:0037
The condition I gave in the \emph{Bulletin de Férussac} for an equation
to be solvable by radicals is too restrictive; there are few exceptions,
but there are some.

% ENSEG:W08:PDF056:0038
The last application of the theory of equations concerns the modular
equations of elliptic functions.

% ENSEG:W08:PDF056:0039
It is known that the group of the equation whose roots are the sines of
the amplitudes of the $p^2-1$ divisions of a period is
% ENSEG:W08:PDF056:0040
\[
x_{k,l},\qquad x_{ak+bl,\;ck+dl};
\]
% ENSEG:W08:PDF056:0041
consequently, the corresponding modular equation will have the group
% ENSEG:W08:PDF056:0042
\[
x_{\frac{k}{l}},\qquad x_{\frac{ak+bl}{ck+dl}},
\]

\GalSourcePageBoundary{057}{L0056}{28}{28}

% ENSEG:W08:PDF057:0043
in which $\dfrac{k}{l}$ can take the $p+1$ values
% ENSEG:W08:PDF057:0044
\[
\infty,\ 0,\ 1,\ 2,\ \ldots,\ p-1.
\]
% ENSEG:W08:PDF057:0045
Thus, agreeing that $k$ may be infinite, one can simply write
% ENSEG:W08:PDF057:0046
\[
x_k,\qquad x_{\frac{ak+b}{ck+d}}.
\]

% ENSEG:W08:PDF057:0047
By giving $a$, $b$, $c$, and $d$ all possible values, one obtains
% ENSEG:W08:PDF057:0048
\[
(p+1)p(p-1)\text{ permutations.}
\]

% ENSEG:W08:PDF057:0049
Now this group decomposes \emph{properly} into two groups whose
substitutions are
% ENSEG:W08:PDF057:0050
\[
x_k,\qquad x_{\frac{ak+b}{ck+d}},
\]
% ENSEG:W08:PDF057:0051
where $ad-bc$ is a quadratic residue of $p$.

% ENSEG:W08:PDF057:0052
The group thus simplified has
% ENSEG:W08:PDF057:0053
\[
(p+1)p\frac{p-1}{2}\text{ permutations.}
\]

% ENSEG:W08:PDF057:0054
But it is easy to see that it is no longer properly decomposable unless
$p=2$ or $p=3$.

% ENSEG:W08:PDF057:0055
Thus, however one transforms the equation, its group will always have
the same number of permutations.

% ENSEG:W08:PDF057:0056
But it is interesting to know whether the degree can be lowered.

% ENSEG:W08:PDF057:0057
First, it cannot be lowered below $p$, since an equation of degree less
than $p$ cannot have $p$ as a factor in the number of permutations of
its group.

% ENSEG:W08:PDF057:0058
Let us therefore see whether the equation of degree $p+1$, whose roots
$x_k$ are designated by giving $k$ all values including infinity, and
whose group has substitutions
% ENSEG:W08:PDF057:0059
\[
x_k,\qquad x_{\frac{ak+b}{ck+d}},
\]
% ENSEG:W08:PDF057:0060
with $ad-bc$ a square, can be lowered to degree $p$.

% ENSEG:W08:PDF057:0061
For this, the group must decompose (improperly, of course) into $p$
groups of $(p+1)\dfrac{p-1}{2}$ permutations each.

\GalSourcePageBoundary{058}{L0057}{29}{29}

% ENSEG:W08:PDF058:0062
Let $0$ and $\infty$ be two associated letters in one of these groups.
The substitutions that do not interchange $0$ and $\infty$ will be of
the form
% ENSEG:W08:PDF058:0063
\[
x_k,\qquad x_{m^2k}.
\]

% ENSEG:W08:PDF058:0064
Thus, if $M$ is the letter associated with $1$, the letter associated
with $m^2$ will be $m^2M$. When $M$ is a square, we shall therefore have
$M^2=1$. But this simplification can occur only for $p=5$.

% ENSEG:W08:PDF058:0065
For $p=7$ one finds a group of $(p+1)\dfrac{p-1}{2}$ permutations in
which
% ENSEG:W08:PDF058:0066
\[
\infty,\ 1,\ 2,\ 4
\]
% ENSEG:W08:PDF058:0067
have respectively as their associated letters
% ENSEG:W08:PDF058:0068
\[
0,\ 3,\ 6,\ 5.
\]
% ENSEG:W08:PDF058:0069
This group has substitutions of the form
% ENSEG:W08:PDF058:0070
\[
x_k,\qquad x_{a\frac{k-b}{k-c}},
\]
% ENSEG:W08:PDF058:0071
where $b$ is the letter associated with $c$, and $a$ is a letter that is
a residue or a nonresidue at the same time as $c$.

% ENSEG:W08:PDF058:0072
For $p=11$, the same substitutions hold with the same notation,
% ENSEG:W08:PDF058:0073
\[
\infty,\ 1,\ 3,\ \GalPrintedError{5}{W08-PE001},\ 5,\ 9
\]
% ENSEG:W08:PDF058:0074
having respectively as their associated letters
% ENSEG:W08:PDF058:0075
\[
0,\ 2,\ 6,\ 8,\ 10,\ 7.
\]

% ENSEG:W08:PDF058:0076
\GalCriticalNote{W08-PE001}{The corrected six-point partition for $p=11$}{The
1897 upper row repeats $5$ and omits the projective point $4$. The corresponding
1846 publication reads $\infty,1,3,4,5,9$. Exact enumeration confirms the
corrected partition
$\{\infty,0\},\{1,2\},\{3,6\},\{4,8\},\{5,10\},\{9,7\}$
and its stabilizer of order $60$ and index $11$ in
$\operatorname{PSL}(2,11)$, whose order is $660$. The 1897 value remains in
the translation, but any use of the $p=11$ construction must apply this
correction.}

% ENSEG:W08:PDF058:0077
Thus, for $p=5$, $7$, and $11$, the modular equation is lowered to
degree $p$.

% ENSEG:W08:PDF058:0078
Strictly speaking, this \GalPrintedError{equation}{W08-PE002} is not
possible in the higher cases.

% ENSEG:W08:PDF058:0079
\GalCriticalNote{W08-PE002}{``Reduction,'' not ``equation''}{The 1897 edition
prints ``equation,'' whereas the 1846 publication prints ``reduction'' and
Tannery's 1908 collation gives the same reading for the manuscript and
Liouville. With that attested correction, the sentence says that this
\emph{reduction} is impossible in higher cases; it does not deny the existence
of the equation itself.}

% ENSEG:W08:PDF058:0080
The third Memoir concerns integrals.

% ENSEG:W08:PDF058:0081
It is known that a sum of terms of the same elliptic function always
reduces to a single term, together with algebraic or logarithmic
quantities.

% ENSEG:W08:PDF058:0082
There are no other functions for which this property holds.

% ENSEG:W08:PDF058:0083
But entirely similar properties take its place for all integrals of
algebraic functions.

\GalSourcePageBoundary{059}{L0058}{30}{30}

% ENSEG:W08:PDF059:0084
One treats simultaneously all integrals whose differential is a
function of the variable and of the same irrational function of the
variable, whether that irrational function is or is not a radical, and
whether it can or cannot be expressed by radicals.

% ENSEG:W08:PDF059:0085
One finds that the number of distinct periods of the most general
integral relative to a given irrational function is always even.

% ENSEG:W08:PDF059:0086
Let this number be $2n$; then one has the following theorem:

% ENSEG:W08:PDF059:0087
Any sum of terms reduces to $n$ terms, together with algebraic and
logarithmic quantities.

% ENSEG:W08:PDF059:0088
Functions of the first kind are those for which the algebraic and
logarithmic part is zero.

% ENSEG:W08:PDF059:0089
There are $n$ distinct such functions.

% ENSEG:W08:PDF059:0090
Functions of the second kind are those for which the complementary part
is purely algebraic.

% ENSEG:W08:PDF059:0091
There are $n$ distinct such functions.

% ENSEG:W08:PDF059:0092
One may suppose that the differentials of the other functions are never
infinite except once, at $x=a$, and, moreover, that their complementary
part reduces to a single logarithm, $\log P$, where $P$ is an algebraic
quantity. Denoting these functions by $\Pi(x,a)$, one has the theorem
% ENSEG:W08:PDF059:0093
\[
\Pi(x,a)-\Pi(a,x)=\sum \varphi a\mathbin{\cdot}\psi x,
\]
% ENSEG:W08:PDF059:0094
where $\varphi a$ and $\psi x$ are functions of the first and second kinds.

% ENSEG:W08:PDF059:0095
Calling $\Pi(a)$ and $\psi$ the periods of $\Pi(x,a)$ and $\psi x$
corresponding to the same circuit of $x$, one deduces
% ENSEG:W08:PDF059:0096
\[
\Pi(a)=\sum\psi\times\varphi a.
\]
% ENSEG:W08:PDF059:0097
Thus the periods of functions of the third kind are always expressed in
terms of functions of the first and second kinds.

% ENSEG:W08:PDF059:0098
One may also deduce theorems analogous to Legendre's theorem
% ENSEG:W08:PDF059:0099
\[
\GalWitnessVariant{FE'+EF'-FF'=\frac{\pi}{2}}{W08-WV001}.
\]

% ENSEG:W08:PDF059:0100
\GalCriticalNote{W08-WV001}{Tannery's manuscript form of Legendre's relation}{The
formula $FE'+EF'-FF'=\pi/2$ printed in 1846 and 1897 is Legendre's standard
relation in the source's $F/K$ notation. A high-resolution reading of
Tannery's 1908 collation shows that the manuscript form begins with $E'$, not
$F'$:
\[
E'F''-E''F'=(\pi/2)\sqrt{-1}.
\]
The two displays express the same relation under an explicit half-period and
quasiperiod normalization. Indeed, writing $F_c=K(k')$ and $E_c=E(k')$, take
$\omega_1=F$, $\omega_2=iF_c$, $\eta_1=E$, and
$\eta_2=i(F_c-E_c)$. Then
$\eta_1\omega_2-\eta_2\omega_1=i(EF_c+E_cF-FF_c)=\pi i/2$.
This corrects the earlier editorial reading of the manuscript formula without
altering the historical text. Nothing later in this volume depends on the
relation.}

% ENSEG:W08:PDF059:0101
The reduction of functions of the third kind to definite integrals,
which is M.~Jacobi's finest discovery, is not practicable outside the
case of elliptic functions.

% ENSEG:W08:PDF059:0102
The multiplication of integral functions by an integer

\GalSourcePageBoundary{060}{L0059}{31}{31}

% ENSEG:W08:PDF060:0103
is always possible, like addition, by means of an equation of degree
$n$ whose roots are the values to be substituted in the integral to
obtain the reduced terms.

% ENSEG:W08:PDF060:0104
The equation that gives the division of the periods into $p$ equal
parts has degree $p^{2n}-1$. Its group has in all
% ENSEG:W08:PDF060:0105
\[
(p^{2n}-1)(p^{2n}-p)\ldots(p^{2n}-p^{2n-1})\text{ permutations.}
\]

% ENSEG:W08:PDF060:0106
The equation that gives the division of a sum of $n$ terms into $p$
equal parts has degree $p^{2n}$. It is solvable by radicals.

% ENSEG:W08:PDF060:0107
\emph{On transformation.} --- First, by following arguments analogous
to those Abel recorded in his last Memoir, one can prove that if, in
the same relation between integrals, one has the two functions
% ENSEG:W08:PDF060:0108
\[
\int \Phi(x,X)\,dx,\qquad \int \Psi(y,Y)\,dy,
\]
% ENSEG:W08:PDF060:0109
the latter integral having $2n$ periods, one may suppose that $y$ and
$Y$ are expressed, by means of a single equation of degree $n$, as
functions of $x$ and $X$.

% ENSEG:W08:PDF060:0110
Accordingly, one may suppose that transformations always take place
between only two integrals, since, on taking any rational function of
$y$ and $Y$, one will evidently have
% ENSEG:W08:PDF060:0111
\[
\sum\!\int f(y,Y)\,dy=\int F(x,X)\,dx+\text{ an alg. and log. quantity.}
\]

% ENSEG:W08:PDF060:0112
There would be evident reductions in this equation if the integrals on
the two sides did not both have the same number of periods.

% ENSEG:W08:PDF060:0113
Thus we need compare only integrals that both have the same number of
periods.

% ENSEG:W08:PDF060:0114
It will be proved that the least degree of irrationality of two such
integrals cannot be greater for one than for the other.

% ENSEG:W08:PDF060:0115
It will then be shown that one can always transform a given integral
into another in which one period of the first is divided by the prime
number $p$, while the other $2n-1$ periods remain the same.

\GalSourcePageBoundary{061}{L0060}{32}{32}

% ENSEG:W08:PDF061:0116
It will therefore remain only to compare integrals whose periods are
the same on both sides and such, consequently, that $n$ terms of one
are expressed, without any equation other than a single one of degree
$n$, in terms of those of the other, and conversely. Here we know
nothing.

\medskip

% ENSEG:W08:PDF061:0117
You know, my dear Auguste, that these subjects are not the only ones I
have explored. For some time my principal meditations have been directed
toward applying the theory of ambiguity to transcendental analysis. The
question was to see \emph{a priori}, in a relation among transcendental
quantities or functions, what interchanges could be made, what quantities
could be substituted for the given quantities, without the relation
ceasing to hold. This immediately reveals the impossibility of many
expressions one might seek. But I have no time, and my ideas are not yet
well developed in this immense field.

% ENSEG:W08:PDF061:0118
You will have this Letter printed in the \emph{Revue encyclopédique}.

% ENSEG:W08:PDF061:0119
I have often ventured in my life to put forward propositions of which I
was not certain; but everything I have written here has been in my mind
for nearly a year, and it is too much in my interest not to be mistaken
for anyone to suspect me of stating theorems of which I did not have a
complete proof.

% ENSEG:W08:PDF061:0120
You will publicly ask Jacobi or Gauss to give their opinion, not on the
truth, but on the importance of the theorems.

% ENSEG:W08:PDF061:0121
After that, I hope there will be people who find it profitable to
decipher all this mess.

% ENSEG:W08:PDF061:0122
I embrace you warmly.

% ENSEG:W08:PDF061:0123
\noindent\hfill\textsc{E. Galois.}\hspace{2.5em}\par
\vspace{0.15em}
% ENSEG:W08:PDF061:0124
\noindent\hspace{2.2em}29 May 1832.\par

\GalArticleEndRule
